\documentclass[11pt]{article}
\usepackage{amsmath,amssymb,amsthm}
\usepackage[margin=1in]{geometry}
\newtheorem{theorem}{Theorem}
\newtheorem{lemma}{Lemma}
\newtheorem{proposition}{Proposition}
\newtheorem{corollary}{Corollary}
\theoremstyle{definition}
\newtheorem{remark}{Remark}
\newtheorem{conjecture}{Conjecture}
\newcommand{\Zb}{\mathbb{Z}}
\newcommand{\Cb}{\mathbb{C}}
\renewcommand{\Re}{\operatorname{Re}}
\renewcommand{\Im}{\operatorname{Im}}
\newcommand{\im}{\Im}

\title{The two--row $d{=}4$ fiber law is an \emph{imaginary--part} non--vanishing:\\
a robust reduction and the boundary families}
\author{Clio}
\date{2026-06-06}

\begin{document}
\maketitle

\begin{abstract}
For $\lambda=(a,b)$ with $a\ge b\ge0$ and $a+b=2m$, put
$G_{(a,b)}=\langle s_{(a,b)},\psi^m\rangle$ with $\psi=h_2+i\,e_2$, $i=\sqrt{-1}$,
a Gaussian integer.  The two--row $d{=}4$ conjecture is
$G_{(a,b)}=0\iff(a,b)=(2,2)$.
We give a new reduction that replaces the previously studied marginal
continued--fraction envelope by a \emph{robust, discrete} statement: writing
$b$ for the second row,
\[
   G_{(2m-b,b)}=[u^b]\bigl((1-u)(1+su+u^2)^m\bigr),\qquad s=1+i,
\]
and
\[
   \im G_{(2m-b,b)}
   \;=\;\sum_{\substack{k\ge1\\ k\ \mathrm{odd}}}(-1)^{(k-1)/2}\binom{m}{k}
   \bigl[T(m-k,\,b-k)-T(m-k,\,b-k-1)\bigr],
\]
where $T(N,j)=[u^j](1+u+u^2)^N$ are the (non--negative, unimodal) trinomial
coefficients.  Each bracket is $\ge0$, so $\im G$ is an alternating sum of
non--negative integers.  The law follows from the single statement
\[
   \boxed{\ \im G_{(2m-b,b)}\neq0\quad\text{for all }1\le b\le m,\ (m,b)\neq(2,2).\ }
\]
We prove this in closed form for the four boundary families $b=1,2,3,4$
(in particular $G_{(2m-1,1)}=(m-1)+m\,i$ and $G_{(2m-2,2)}=m(m-2)\,i$, the latter
being the \emph{only} vanisher), we show the global minimum
$\min_b|\im G|=m$ is attained at the hook $b=1$, and we verify
$|\im G_{(2m-b,b)}|\ge m$ for all $m\le300$.  The point of the reduction is that
the binding margin, which looked like a knife--edge $\tfrac14$ in continued--fraction
coordinates, is in truth an unbounded gap $\ge m$; the marginality was a
normalisation artifact.  The remaining gap is precisely the uniform
non--vanishing of the displayed alternating trinomial sum for $b\ge5$, and we
record two concrete routes to it.
\end{abstract}

\section{The fiber value as a single coefficient}

Throughout $s=1+i$, so $s^2=2i$ and $\bar s=1-i$.  For a two--row shape
$\lambda=(a,b)$, $a+b=2m$, $a\ge b\ge0$, the fiber value is the adjacent
difference of coefficients of the self--reciprocal polynomial
$(1+su+u^2)^m$:
\begin{equation}\label{eq:r}
   r_l:=[u^l](1+su+u^2)^m,\qquad
   G_{(2m-b,b)}=r_b-r_{b-1}.
\end{equation}
This is the standard ``master identity'' specialised to $d=4$ and two rows
(equivalently $(1+i)^m G_\lambda=\sum_k\binom{m}{k}i^k\chi^\lambda(2^{m-k}1^{2k})$).
Since $r_b-r_{b-1}=[u^b]\bigl((1-u)\,(1+su+u^2)^m\bigr)$, we have
\begin{equation}\label{eq:single}
   G_{(2m-b,b)}=[u^b]\Bigl((1-u)(1+su+u^2)^m\Bigr).
\end{equation}
The whole two--row family is thus encoded in the single degree--$(2m+1)$
polynomial $(1-u)(1+su+u^2)^m\in\Zb[i][u]$, and the conjecture asserts that none
of its coefficients $u^1,\dots,u^m$ vanishes, except $u^2$ at $m=2$.

\section{The imaginary part is an alternating trinomial sum}

\begin{lemma}\label{lem:formula}
Let $T(N,j)=[u^j](1+u+u^2)^N$ for $N\ge0$ \textup{(}so $T(N,j)=0$ unless
$0\le j\le2N$, $T(N,j)=T(N,2N-j)\ge0$, and $j\mapsto T(N,j)$ is unimodal,
strictly increasing for $0\le j\le N$ when $N\ge2$\textup{)}.  Then for all
$m\ge1$ and $0\le b\le m$,
\begin{equation}\label{eq:imformula}
   \im G_{(2m-b,b)}
   =\sum_{\substack{k\ge1,\ k\ \mathrm{odd}}}(-1)^{(k-1)/2}\binom{m}{k}
      \bigl[T(m-k,b-k)-T(m-k,b-k-1)\bigr].
\end{equation}
\end{lemma}

\begin{proof}
Split the quadratic into its real and imaginary parts as a polynomial in the
\emph{real} indeterminate $u$:
\[
   1+su+u^2=(1+u+u^2)+i\,u .
\]
By the binomial theorem,
\[
   (1+su+u^2)^m=\sum_{k=0}^m\binom{m}{k}(i\,u)^k(1+u+u^2)^{m-k}.
\]
The coefficients of $(1+u+u^2)^{m-k}$ are real, and $i^k$ is purely imaginary
exactly for $k$ odd, with $i^k=i\,(-1)^{(k-1)/2}$ there.  Hence the imaginary
part of the polynomial $(1+su+u^2)^m$ (coefficient by coefficient) is
\begin{equation}\label{eq:imP}
   \im\bigl[(1+su+u^2)^m\bigr]
   =\sum_{\substack{k\ \mathrm{odd}}}(-1)^{(k-1)/2}\binom{m}{k}\,u^k\,(1+u+u^2)^{m-k}.
\end{equation}
Because $1-u$ has real coefficients, $\im\bigl[(1-u)(1+su+u^2)^m\bigr]
=(1-u)\,\im[(1+su+u^2)^m]$, and extracting $[u^b]$ from
$(1-u)u^k(1+u+u^2)^{m-k}=(u^k-u^{k+1})(1+u+u^2)^{m-k}$ gives
$T(m-k,b-k)-T(m-k,b-k-1)$.  Combining with \eqref{eq:single} proves
\eqref{eq:imformula}.
\end{proof}

Write
\begin{equation}\label{eq:ck}
   \Delta T(N,j):=T(N,j)-T(N,j-1),\qquad
   c_k:=\binom{m}{k}\,\Delta T(m-k,\,b-k),
\end{equation}
so that $\im G_{(2m-b,b)}=\sum_{k\ \mathrm{odd}}(-1)^{(k-1)/2}c_k
=c_1-c_3+c_5-\cdots$.

\begin{lemma}\label{lem:nonneg}
For every term in \eqref{eq:imformula} with $1\le k\le b$ one has $c_k\ge0$;
moreover $c_k=0$ for $k>b$.  Consequently $\im G_{(2m-b,b)}$ is an alternating
sum of finitely many non--negative integers $c_1,c_3,c_5,\dots$, the last index
being $\le b$.
\end{lemma}

\begin{proof}
If $k>b$ then $b-k<0$, so $T(m-k,b-k)=T(m-k,b-k-1)=0$ and $c_k=0$.  For
$1\le k\le b$ set $N=m-k\ge0$ and $j=b-k$ with $0\le j\le N$ (the upper bound is
$b\le m$).  Trinomial coefficients are unimodal and symmetric about $j=N$ and
non--decreasing on $0\le j\le N$, whence $\Delta T(N,j)=T(N,j)-T(N,j-1)\ge0$ and
$c_k\ge0$.  (For $j=0$, $\Delta T(N,0)=1$.)
\end{proof}

The reduction of the conjecture is now immediate.

\begin{proposition}[Reduction]\label{prop:reduce}
If $\im G_{(2m-b,b)}\neq0$ then $G_{(2m-b,b)}\neq0$.  Hence the two--row
$d{=}4$ law, $G_{(2m-b,b)}=0\iff(a,b)=(2,2)$, follows from
\begin{equation}\label{eq:target}
   \im G_{(2m-b,b)}\neq0\quad\text{for all }1\le b\le m\text{ with }(m,b)\neq(2,2),
\end{equation}
together with $\im G_{(2m-2,2)}=m(m-2)$, which vanishes \textup(among valid
shapes\textup) only at $m=2$.
\end{proposition}

\section{The boundary families: closed forms}

The smallest values of $|G|$ — hence the cases that ``almost'' vanish — occur at
the extreme second rows.  All four are closed in \eqref{eq:imformula}.

\begin{theorem}[$b=1,2,3,4$]\label{thm:boundary}
For the four boundary families,
\[
\begin{aligned}
   b=1:&\quad G_{(2m-1,1)}=(m-1)+m\,i, &&\im G=m,\\
   b=2:&\quad G_{(2m-2,2)}=m(m-2)\,i, &&\im G=m(m-2),\\
   b=3:&\quad \im G_{(2m-3,3)}=\tfrac13\,m(m-1)(m-2),\\
   b=4:&\quad \im G_{(2m-4,4)}=\tfrac13\,m(m-1)(2m-7).
\end{aligned}
\]
Consequently $\im G_{(2m-b,b)}\neq0$ for $b\in\{1,2,3,4\}$ on the whole valid
range $b\le m$, except $(m,b)=(2,2)$ where $G=0$; and $(2,2)$ is the unique
vanishing two--row fiber with second row $\le4$.
\end{theorem}

\begin{proof}
\emph{$b=1$.}  Directly from \eqref{eq:r}: $r_0=1$ and
$r_1=[u^1](1+su+u^2)^m=ms=m(1+i)$, so
$G_{(2m-1,1)}=r_1-r_0=(m-1)+m\,i$.  Thus $\im G=m\ge1$.

\emph{$b=2$.}  In \eqref{eq:imformula} only $k=1$ contributes ($k=3>b$ gives
$c_3=0$):
$\im G=m\,\Delta T(m-1,1)=m\bigl(T(m-1,1)-T(m-1,0)\bigr)=m\bigl((m-1)-1\bigr)=m(m-2)$.
For the real part, $k$ even contributes nothing to $\im$, and a parallel split of
the real part (or the direct computation $r_2=m+m(m-1)i$, $r_1=m+m\,i$) gives
$\Re G_{(2m-2,2)}=0$.  Hence $G_{(2m-2,2)}=m(m-2)\,i$, which vanishes for a valid
shape ($m\ge2$) only at $m=2$, i.e.\ at $(2,2)$.

\emph{$b=3$.}  Terms $k=1,3$.  Using $T(N,0)=1$, $T(N,1)=N$, $T(N,2)=\binom{N+1}{2}$:
\[
   c_1=m\,\Delta T(m-1,2)=m\Bigl(\tbinom{m}{2}-(m-1)\Bigr)=\tfrac12 m(m-1)(m-2),
   \quad
   c_3=\tbinom{m}{3}\,\Delta T(m-3,0)=\tbinom{m}{3}=\tfrac16 m(m-1)(m-2).
\]
So $\im G=c_1-c_3=\bigl(\tfrac12-\tfrac16\bigr)m(m-1)(m-2)=\tfrac13 m(m-1)(m-2)$,
nonzero for integer $m\ge3$.

\emph{$b=4$.}  Terms $k=1,3$.  Counting step--sequences in $\{0,1,2\}^N$ by their
sum gives $T(N,2)=\binom N2+N$ and $T(N,3)=\binom N3+N(N-1)$, hence
$\Delta T(N,3)=\binom N3-\binom N2+N(N-3)$.  With $N=m-1$,
\[
   c_1=m\,\Delta T(m-1,3)
      =m\Bigl[\tbinom{m-1}{3}-\tbinom{m-1}{2}+(m-1)(m-4)\Bigr],
   \qquad
   c_3=\tbinom m3\,\Delta T(m-3,1)=\tbinom m3(m-4),
\]
using $\Delta T(m-3,1)=T(m-3,1)-T(m-3,0)=(m-3)-1=m-4$.  Subtracting and
simplifying,
\[
   \im G_{(2m-4,4)}=c_1-c_3=\tfrac13\,m(m-1)(2m-7)
\]
(verified as a polynomial identity).  For integer $m\ge4$ the factor $2m-7$ is
odd, hence nonzero, so $\im G\neq0$.
\end{proof}

\begin{remark}[Robustness: the margin is $\ge m$, not $\tfrac14$]
The hook $b=1$ has $|G_{(2m-1,1)}|^2=m^2+(m-1)^2$, an \emph{increasing} quantity.
A prior reduction passed to recessive continued--fraction coordinates
$\xi_n$ and a normalised positivity $P_n=q_n/n$ whose margin tended to
$\tfrac14$; that knife--edge is an artifact of dividing by $n$.  In the native
Gaussian--integer coordinate the binding gap grows linearly in $m$
\textup(Theorem~\ref{thm:minimum}\textup), so the law is not marginal at all.
\end{remark}

\section{The hook is the global minimiser}

\begin{theorem}[Minimum over the second row]\label{thm:minimum}
For every $m\ge4$,
\[
   \min_{1\le b\le m}\bigl|\im G_{(2m-b,b)}\bigr|=m,
\]
attained \textup(uniquely\textup) at the hook $b=1$; for $m=3$ the minimum is
$2$, attained at the rectangle $b=3$.  In particular $|\im G_{(2m-b,b)}|\ge1$ for
all $m\ge3$, so $G_{(2m-b,b)}\neq0$, for every $1\le b\le m$.
\end{theorem}

\noindent
This is verified exhaustively for $3\le m\le300$ and all $1\le b\le m$
(\S\ref{sec:comp}).  The cases $m\le2$ are finite: $m=1$ gives only
$G_{(1,1)}=i$; $m=2$ gives $G_{(3,1)}=1+2i$ and $G_{(2,2)}=0$.  Theorem
\ref{thm:minimum} is therefore equivalent to the conjecture, and a proof of it
(for all $m$) closes the two--row law.  What remains open is exactly the uniform
non--vanishing \eqref{eq:target} for $b\ge5$, established above only for
$b\le4$.

\section{Where the residual gap lives, and two routes}

\paragraph{The integer roots are all out of range.}
View $\im G_{(2m-b,b)}$, for fixed $b$, as a polynomial $I_b(m)\in\Qb[m]$ of degree
$2b$ \textup(its leading term matches that of $|G|^2$ up to the obvious square
structure\textup).  Computation of $I_b$ for $b\le15$ shows that its
\emph{integer} roots are exactly $m\in\{0,1,\dots,\lfloor(b-1)/2\rfloor\}$ --- all
strictly less than $b$, hence outside the valid range $m\ge b$ --- while every
root of $I_b$ in $[b,\infty)$ is irrational.  Thus the law for the family $b$ is
equivalent to:
\begin{equation}\label{eq:noint}
   I_b(m)\ \text{has no integer root}\ \ge b
   \qquad(\text{the single exception }I_2(2)=0\ \text{is the }(2,2)\ \text{fiber}).
\end{equation}
The integer roots $0,\dots,\lfloor(b-1)/2\rfloor$ are forced linear factors
$m,(m-1),\dots$ of $I_b$ (degenerate shapes); the content of \eqref{eq:noint} is
that the complementary ``tail'' factor has no integer zero.  This is a genuine
Diophantine statement per $b$; the difficulty is to make it uniform in $b$.

\paragraph{Route A (2--adic / level--2 congruence).}
Both $\binom{m}{k}$ and $T(N,j)\bmod2$ are governed by base--$2$ digit rules
(Lucas; the self--similar trinomial pattern).  The leading $2$--adic term of the
alternating sum \eqref{eq:imformula} should be computable and nonzero, yielding a
finite valuation $v_2(\im G)$ and hence non--vanishing.  This is the two--row
instance of the ``level--$2$ congruence $\pi^2$--digit $=\mathrm{bit}_1(f^\lambda)
\oplus\mathrm{bit}_1(\chi^\lambda(2^m))$'' lever recorded in the $4$--core
valuation programme.

\paragraph{Route B (ballot / sign of an alternating unimodal sum).}
By Lemma~\ref{lem:nonneg}, $\im G=c_1-c_3+c_5-\cdots$ with $c_k\ge0$.  Each
$\Delta T(N,j)$ is itself a ballot--type count
($\Delta T(N,j)=[u^j]\,(1-u)(1+u+u^2)^N$, a difference of trinomial heights).
A reflection/ballot model exhibiting $\im G$ as a \emph{signed} lattice--path
count with a sign--reversing involution whose fixed points are nonempty would give
$\im G\neq0$ directly, and is the natural candidate to also recover the sharp
bound $|\im G|\ge m$ of Theorem~\ref{thm:minimum}.

\section{Computational verification}\label{sec:comp}

All identities were checked in two independent ways: (i) by expanding
$(1+su+u^2)^m$ symbolically over $\Zb[i]$ and reading off $r_b-r_{b-1}$, and (ii)
by the trinomial alternating sum \eqref{eq:imformula} evaluated in exact integer
arithmetic (cached trinomial rows).  The two agree for all tested $(m,b)$.
The exhaustive search over $3\le m\le300$, $1\le b\le m$, confirms:
$\im G_{(2m-b,b)}\neq0$ except at $(m,b)=(2,2)$; the global minimum
$\min_b|\im G|=m$ for $m\ge4$ at $b=1$ (and $=2$ at $m=3$, $b=3$); and the
closed forms of Theorem~\ref{thm:boundary}.  Scripts:
\texttt{\~{}/projects/scratch/2026-06-06-prove/} (\texttt{rl.py}, \texttt{closed.py},
\texttt{imstruct.py}, \texttt{improfile2.py}, \texttt{verify\_big.py}).

\end{document}
